\documentclass[11pt]{article}
\usepackage{color}
\usepackage{amsmath,amsthm,amssymb,multirow,paralist}
\usepackage[margin=0.8in]{geometry}
\usepackage{hyperref}

\begin{document}

\begin{center}
{\Large \textbf{COM S 4740/5740: Introduction to Machine Learning}\\Homework \#3}\\

\linethickness{1mm}\line(1,0){498}

\begin{enumerate}
\item Please put required code files and report into a
compressed file ``HW\#\_FirstName\_LastName.zip''
\item Unlimited number of submissions are
allowed on Canvas and the latest one will be graded.
\item {\color{red} No later submission is accepted.}
\item Please read and follow submission instructions. No exception
will be made to accommodate incorrectly submitted files/reports.
\item All students are required to typeset their reports using
latex. Overleaf
(\url{https://www.overleaf.com/learn/latex/Tutorials}) can be a
good start.
\end{enumerate}

\linethickness{1mm}\line(1,0){498}

\end{center}

%%%%%%%%%%%%%%%%%%%%%%%%%%%%%%%%%%%%%%%%%%%%%%%%%%%%%%%%%%%%%%%%%%%%%%%%%%%%%%%

%%%%%%%%%%%%%%%%%%%%%%%%%%%%%%%%%%%%%%%%%%%%%%%%%%%%%%%%%%%%%%%%%%%%%%%%%%%%%%%


\begin{enumerate}

\item (20 points) Consider the toy data set $\{([0, 0], -1),
([2, 2], -1), ([2, 0], +1)\}$. Set up the dual problem for
the toy data set. Then, solve the dual problem and compute
$\alpha^*$, the optimal Lagrange multipliers. (Note that there will
be three weights $\boldsymbol w = [w_0, w_1, w_2]$ by considering
the bias.)

{\color{red}

Augmenting w0: \\
$x_1 = [1, 0, 0], y_1 = -1$ \\
$x_2 = [1, 2, 2], y_2 = -1$ \\
$x_3 = [1, 2, 0], y_3 = 1$

The dual problem:
\[ \max_{\alpha} \sum_{i=1}^3 \alpha_i - \frac{1}{2} \sum_{i=1}^3 \sum_{j=1}^3 \alpha_i \alpha_j y_i y_j (x_i \cdot x_j) \]

Calculating $x_ix_j $ for each value i and j from the sums:\\
$x_1 \cdot x_1 = 1, x_1 \cdot x_2 = 1, x_1 \cdot x_3 = 1$ \\
$x_2 \cdot x_2 = 9, x_2 \cdot x_3 = 5, x_3 \cdot x_3 = 5$

Plugging in:
\[ L(\alpha) = (\alpha_1 + \alpha_2 + \alpha_3) - \frac{1}{2}(\alpha_1^2 + 9\alpha_2^2 + 5\alpha_3^2 + 2\alpha_1\alpha_2 - 2\alpha_1\alpha_3 - 10\alpha_2\alpha_3) \]

Taking partial derivatives with respect to each $\alpha_i$ and setting to zero:
\begin{enumerate}
    \item $\frac{\partial L_D}{\partial \alpha_1} = 1 - \alpha_1 - \alpha_2 + \alpha_3 = 0$
    \item $\frac{\partial L_D}{\partial \alpha_2} = 1 - \alpha_1 - 9\alpha_2 + 5\alpha_3 = 0$
    \item $\frac{\partial L_D}{\partial \alpha_3} = 1 + \alpha_1 + 5\alpha_2 - 5\alpha_3 = 0$
\end{enumerate}

Solving the system of equations using elimination:\\
eliminate $\alpha_1$ and $\alpha_3$:
    \begin{align*}
    (1 - \alpha_1 - 9\alpha_2 + 5\alpha_3) & = 0 \\
    + (1 + \alpha_1 + 5\alpha_2 - 5\alpha_3) & = 0 \\
     \hline
    2 - 4\alpha_2 & = 0 \implies \alpha_2^* = 0.5
    \end{align*}

eliminate $\alpha_1$:
    \begin{align*}
    (1 - \alpha_1 - \alpha_2 + \alpha_3) & = 0 \\
    + (1 + \alpha_1 + 5\alpha_2 - 5\alpha_3) & = 0 \\
     \hline
    2 + 4\alpha_2 - 4\alpha_3 & = 0
    \end{align*}


substitute $\alpha_2 = .5$ into eliminated equation:
$2 + 4(.5) - 4(\alpha_3) = 0 \implies 4 - 4\alpha_3 = 0 \implies \alpha_3 = 1$

substitute $\alpha_2 = .5 , \alpha_3 =1$ into equation (a):
$1 - \alpha_1 - .5 + 1 = 0 \implies 1.5 = \alpha_1 $

The optimal Lagrange multipliers are:
\[ \alpha^* = [1.5, 0.5, 1.0] \]

}

\item (20 points) In a separable case, when a multiplier
$\alpha_i > 0$, its corresponding data point $(\boldsymbol x_i,
y_i)$ is on the boundary of the optimal separating hyperplane
with $y_i(\boldsymbol w^T \boldsymbol x_i) = 1$.

Show that the inverse is not True. Namely, it is possible that
$\alpha_i = 0$ and $(\boldsymbol x_i, y_i)$ is on the boundary
satisfying $y_i(\boldsymbol w^T \boldsymbol x_i) = 1$.

[Hint: Consider a toy data set with two positive examples at
([0,0], +1) and ([1, 0], +1), and one negative example at ([0,
1], -1).] (Note that there will be three weights $\boldsymbol w =
[w_0, w_1, w_2]$ by considering the bias.)


{\color{red}
 $y_i(\boldsymbol{w}^T \boldsymbol{x}_i) \ge 1$ 
\begin{enumerate}
    \item $1(w_0 + w_1(0) + w_2(0)) \ge 1 \implies w_0 \ge 1$
    \item $1(w_0 + w_1(1) + w_2(0)) \ge 1 \implies w_1 + w_0 \ge 1$
    \item $-1(w_0 + w_1(0) + w_2(1)) \ge 1 \implies -(w_0+w_2) \ge 1 \implies w_0 + w_2\le -1$
\end{enumerate}

Minimizing $||\boldsymbol{w}||^2 \implies \min(w_0^2 + w_1^2 + w_2^2)$
We need to find the smallest value each weight can be

$w_0: w_0 \ge \textbf{1}$ the smallest value it can be is 1

$w_1:  w_1 + 1 \ge 1 \implies w_1 \ge \textbf{0}$ the smallest value it can be is 0

$w_2: 1 + w_2\le -1 \implies w_2 \le -2 $ the smallest value it can be is -2

$\boldsymbol{w}^* = [1,0,-2]$\\

To be on the margin boundary, if $y_i(\boldsymbol{w}^T \boldsymbol{x}_i) = 1$\\
point 1: $1(1 + 0(0) + -2(0)) = 1$ \checkmark \\
point 2: $1(1 + 0(1) + -2(0)) = 1$ \checkmark \\
point 3: $-1(1 + 0(0) + -2(1)) = 1$ \checkmark \\

$$ \boldsymbol{w} = \sum_{i=1}^{3} \alpha_i y_i \boldsymbol{x}_i $$

$[1,0,-2] = \alpha_1(1)[1,0,0] + \alpha_2(1)[1,1,0] + \alpha_3(-1)[1,0,1]$ \\
$[1,0,-2] = [\alpha_1,0,0] + [\alpha_2,\alpha_2,0] + [-\alpha_3,0,-\alpha_3]$\\
Taking the 2nd term: $0 = 0+\alpha_2 + 0 \implies \alpha_2 = 0$

Since we showed $\alpha_2 = 0$ and that $(x_2,y_2)$ is on the hyperplane, we have proved that the inverse is not true. 
}


\item (20 points) \textbf{Kernel Function:} A function $K$
computes $K(\boldsymbol x_i, \boldsymbol x_j) =
-\boldsymbol{x_i}^T \boldsymbol x_j $. Is this function a valid
kernel function for SVM? Prove or disprove it.

{\color{red}{

To be a valid kernel function, the function must be both symmetric and positive semi-definite. \\

The function is symmetric because the standard vector dot product is commutative. For any vectors $\boldsymbol{x}_i$ and $\boldsymbol{x}_j$:
\[ K(\boldsymbol{x}_i, \boldsymbol{x}_j) = -\boldsymbol{x}_i^T \boldsymbol{x}_j = -\boldsymbol{x}_j^T \boldsymbol{x}_i = K(\boldsymbol{x}_j, \boldsymbol{x}_i) \]
The dot product is the same in both cases \\


Checking for PSD: \\

A matrix is PSD if for any real column vector u, it holds that: $u^TKu \ge 0$ \\
(I got this from the professor) It is also found here: \hyperlink{}{http://theanalysisofdata.com/probability/C_4.html} \\

For $K(x_ix_j) = -x_i^Tx_j$, the kernel matrix is -X

$ u^T(-x^T_ix_j)u \ge 0$  \\
$ u^T(-X)u \ge 0 $ \\
$ -(u^TXu) \ge 0$ \\

Since we established that $u^TXu \ge 0$, we know that $ -(u^TXu) \le 0$ \\
Because we showed that  $u^TKu \ge 0$ is not true, it is not PSD.

We showed that for this kernel function it is symmetric, but not PSD. Therefore, it is not a valid kernel function for SVM.
}


\end{enumerate}

\end{document}